% ============================================================
% chapter_01_introduction.tex - Chapter 1, Introduction
%
% Purpose: motivates the thesis (satellite-precipitation biases over the
%          Indonesian Maritime Continent, the verification-decomposition
%          gap, the reproducibility gap) and states the research question.
% Sections: 1.1 Background and Motivation, 1.2 Problem Statement,
%           1.3 Research Question, 1.4 Objectives, 1.5 Scope and Limits,
%           1.6 Thesis Outline.
% Included by main_thesis.tex.
% License: CC BY 4.0 (see ../LICENSE).
% ============================================================
%   - paper/additional_information.md (§1 scope, §7.7 limitations)
%   - paper/thesis_outline.md (§1.1-§1.6 mapping)
%
% Conventions: no em-dashes; no AI-language; no overclaim;
%              dual contribution (scientific + reproducibility) seeded here;
%              sensitivity parameters acknowledged as not formally optimised.
% ============================================================

\chapter{Introduction}\label{chap:introduction}

% ============================================================
\section{Background}\label{sec:intro-background}
% ============================================================

Daily-scale precipitation estimates underpin flash-flood early
warning, reservoir release scheduling, drought monitoring, crop
water budgeting, and disaster risk assessment. For most of the
global land surface, the only source of spatially continuous
daily precipitation is the satellite family of products. The
Integrated Multi-satellitE Retrievals for the Global
Precipitation Measurement mission (IMERG), in its version 7
release, has become the de facto reference for this class of
data \citep{ref:hou2014, ref:huffman2020}. Among its variants,
the Late Run (IMERG-L) is the one suited to near-real-time use:
its latency of approximately 14 hours fits within the same-day
decision window of operational flood and drought services, where
the gauge-merged Final Run, released months later, is not yet
available. Because IMERG-L relies on satellite observations
without gauge assimilation, it also provides an uncontaminated
baseline against which a bias correction can be evaluated.

IMERG-L carries systematic biases against gauge observations.
Three modes recur across validation studies: over-detection of
light rain, in which the satellite reports drizzle on days the
gauge records as dry; under-detection of heavy events, in which
the satellite truncates the upper tail that drives flood
thresholds; and a full distributional mismatch in which the
shape of the daily rainfall distribution departs from the gauge
distribution \citep{ref:maggioni2016, ref:tang2020}. These biases
are consequential everywhere, but they are acute over the
tropical Maritime Continent, where convective rainfall dominates
the seasonal cycle and the contributing gauge network is dense
only over Java and southern Sumatra, thinning sharply toward
Maluku and Papua \citep{ref:ramadhan2022, ref:chen2008}.

Four families of statistical and machine-learning methods have
been developed to correct these biases, each addressing a
distinct dimension. Linear Scaling corrects the mean through a
single multiplicative factor \citep{ref:teutschbein2012};
Empirical Quantile Mapping with a parametric fit aligns the full
marginal distribution \citep{ref:gudmundsson2012, ref:themessl2012};
Generalized Pareto Distribution tail adjustment grafts an
extreme-value model onto the upper tail \citep{ref:coles2001};
and Convolutional Neural Networks refine the residual spatial
structure of extreme pixels \citep{ref:pan2019, ref:banomedina2020}.
A hybrid pipeline that integrates all four in sequence, denoted
LSEQM+DL (Linear Scaling (LS), Empirical Quantile Mapping (EQM), 
Generalized Pareto Distribution (GPD), and Convolutional Neural Network 
(CNN) refinement, a Deep Learning (DL) architecture), has been applied 
to IMERG-L over Indonesia in a parallel manuscript \citep{ref:istanto2026mdpi}\MDPIstatus 
and is the subject of this thesis. The detailed review of these method
families and their known limits is given in
Chapter~\ref{chap:litreview}.

Two gaps in the bias correction literature motivate the present
study. The first is a verification gap. Most bias correction
studies report an aggregate skill improvement and stop there,
without attributing the movement of individual metrics to the
specific correction stage responsible, and without testing the
result against gauges that are independent of the correction
reference. As a consequence, it is often unclear which bias
dimension a given pipeline actually corrects, and whether the
reported skill survives a comparison against truly out-of-sample
observations. The second is a reproducibility gap. The
reproducibility discourse in computational science and the FAIR
(Findable, Accessible, Interoperable, Reusable) principles for
scientific data have made independent
verifiability a methodological expectation
\citep{ref:peng2011, ref:stodden2014, ref:wilkinson2016}, yet
the bias correction literature has adopted it unevenly: many
influential studies publish a method without releasing the code
or data, leaving any claim of reproduction to a reader who must
re-implement the method from the paper alone. This thesis is
organised around closing both gaps at once.


% ============================================================
\section{Problem Statement}\label{sec:intro-problem}
% ============================================================

The two gaps combine into a single problem with two faces. On
the scientific side, the question is not whether a hybrid bias
correction pipeline improves daily satellite precipitation in
aggregate, which the literature already establishes, but which
bias dimension each stage corrects, and how far the correction
can reach along each dimension when judged against independent
gauges. A pipeline can report a high aggregate score while
leaving a specific operationally important dimension unmoved, and
only a stage-resolved verification against out-of-sample stations
can reveal this. The day-by-day temporal pairing between
satellite and gauge is the dimension most at risk of being
obscured by an aggregate score, because it is bounded by the raw
retrieval and cannot be lifted by a marginal correction, yet it
is the dimension that timing-sensitive applications depend on.

On the methodological side, the problem is that a verification
result of this kind is only credible if an independent reader can
reconstruct it. A claim that a particular stage moves a
particular metric, or that a particular metric cannot be moved,
is a claim about a specific pipeline applied to specific data,
and it is verifiable only if the pipeline and its configuration
are available to re-run. The problem this thesis addresses is
therefore to deliver a stage-resolved, out-of-sample verification
of a hybrid bias correction pipeline, and to deliver it as a
reproducible artefact that an independent reader can re-execute.


% ============================================================
\section{Objectives and Benefits}\label{sec:intro-objectives}
% ============================================================

The study pursues four objectives:

\begin{enumerate}[noitemsep, topsep=2pt]
    \item To attribute the skill of the LSEQM+DL pipeline
    stage-by-stage against 171 independent Indonesian Meteorological, Climatological, and Geophysical Agency (BMKG) stations,
    establishing which bias dimension each of the four stages
    corrects and by how much.
    \item To test, empirically and theoretically, the bound on
    the Pearson correlation between corrected daily satellite and
    gauge values under marginal correction, and to locate the
    observed value within the empirical band reported for daily
    IMERG over the tropics.
    \item To deliver the pipeline as a reproducible, open,
    configuration-driven implementation that an independent reader
    can re-execute end-to-end on commodity hardware.
    \item To document the operational implications of the result,
    distinguishing the application classes the corrected product
    serves well from those it does not.
\end{enumerate}

The benefits follow the objectives. The scientific benefit is
attribution clarity: a verification posture that separates a
structural limit of a method from a failure of a method, so that
the corrected product is neither over-trusted nor dismissed. The
operational benefit is guidance: a clear statement of which
applications the corrected daily product supports and which it
does not. The open-science benefit is independent verifiability:
a released pipeline that lets any reader reproduce the result,
extend it to another region or product, or challenge it on its
own terms.


% ============================================================
\section{Scope, Limitations, and Constraint}\label{sec:intro-scope}
% ============================================================

\medskip\noindent\textbf{Scope.}
The study corrects IMERG-L over the Indonesian domain
($95$ to $141^{\circ}$E, $11^{\circ}$S to $6^{\circ}$N) at the
native $0.1^{\circ}$ daily resolution, calibrated against the
Climate Prediction Center (CPC) Unified Gauge-Based Analysis of Daily
Precipitation (CPC-UNI) gauge analysis and validated against 171 BMKG stations.
The temporal unit is the dekad (the three 10-day periods within
each month), and statistics are pooled across the 36 dekadal
windows of the 2001 to 2025 record. The four-stage LSEQM+DL
pipeline and its station-density-aware blending are taken as the
object of study; the contribution is the verification and the
reproducible implementation, not a new correction algorithm.

\medskip\noindent\textbf{Limitations.}
Three parameters in the pipeline were set by domain reasoning
rather than by a formal optimisation: the CNN blending weight
($\alpha = 0.70$), the GPD threshold (the $80$th wet-day
percentile), and the station-density saturation count
($c_{\text{sat}} = 2$). Their values are defensible and their
direction of effect is understood, but a formal grid search with
a strict temporal hold-out has not been performed and is
identified as future work. The validation is value-independent
but not strictly location-disjoint: the BMKG station values do
not enter the correction, but BMKG station locations contribute
to the density mask through the CPC-UNI network.

\medskip\noindent\textbf{Constraint.}
The strongest constraint on the attainable result is the quality
of the reference. The pipeline calibrates against CPC-UNI, which
is itself a gauge-interpolated analysis with known undercatch of
heavy rainfall in sparse-gauge regions and smoothing at its
$0.5^{\circ}$ resolution. The corrected product cannot be better
than the reference it is calibrated against, except where the
independent BMKG comparison reveals that the correction has moved
the product closer to the true gauge distribution than CPC-UNI
itself. This constraint, and the upstream timing accuracy of the
raw retrieval, bound the result in ways no post-processing can
remove.


% ============================================================
\section{Novelty}\label{sec:intro-novelty}
% ============================================================

The novelty of the thesis is dual. The scientific novelty is the
explicit, stage-resolved demonstration, both empirical and
theoretical, that the Pearson correlation between corrected daily
satellite and gauge values is bounded by the raw satellite timing
rather than by the correction method, paired with the
demonstration that the same pipeline moves every marginal-distribution
metric to the gauge target. Reporting the unchanged correlation
alongside the corrected distribution, and explaining the contrast
as structural rather than as failure, is the verification
contribution. The methodological novelty is the delivery of the
entire result as a reproducible, open, end-to-end pipeline that
an independent reader can re-execute on free cloud infrastructure,
which closes the reproducibility gap for this class of study and
makes the scientific claim independently checkable.


% ============================================================
\section{Hypotheses}\label{sec:intro-hypotheses}
% ============================================================

The study tests two hypotheses:

\begin{description}[noitemsep, topsep=2pt]
    \item[H1.] Each correction stage moves the verification
    metric matched to its design dimension toward the gauge
    target: Linear Scaling moves the mean, Empirical Quantile
    Mapping moves the distribution shape, the Generalized Pareto
    adjustment moves the upper tail, and the Convolutional Neural
    Network refines the spatial structure of the extremes.
    \item[H2.] The Pearson correlation between corrected daily
    satellite and gauge values does not improve across the
    correction stages by a margin exceeding the single-pixel
    sampling uncertainty, because marginal correction preserves
    the rank order of the satellite series and is bounded by the
    day-by-day pairing quality of the raw retrieval.
\end{description}


% ============================================================
\section{Thesis Structure}\label{sec:intro-structure}
% ============================================================

The remainder of the thesis is organised as follows.
Chapter~\ref{chap:litreview} reviews the satellite precipitation
products, the bias dimensions and verification standards, the
bias correction method families, the theoretical bounds on
marginal correction, the reproducible-research literature, and
the Maritime Continent context.
Chapter~\ref{chap:methods} describes the study area and data, the
four-stage LSEQM+DL pipeline, the station-density-aware blending,
the reproducible implementation, the validation framework, and
the verification metric catalogue.
Chapter~\ref{chap:results} reports the stage-wise headline skill,
the stage attribution, the threshold-dependent detection, the
correlation result, the spatial skill distribution, the
per-station Taylor analysis, the temporal stability, the
sensitivity discussion, and the reproducibility evaluation.
Chapter~\ref{chap:discussion} develops the theoretical bound on
the correlation, the operational implications, the
reproducibility implications, the paths forward, and the
limitations.
Chapter~\ref{chap:conclusion} returns to the title question and
states the recommendations.
